\section{Connected activity on the first-visit tree}\label{sec:thinning}

We use the arithmetic estimates of \Cref{sec:paths} to simplify the
closed lines contributing to the high trace.  The restrictions will
depend only on the numerical line and a specified set of residue
coordinates.  They will leave the residue of every center prime
occurring only once on the line free for the later signed average.

\subsection{The first-visit tree and active vertices}

For a closed line $\boldsymbol d$, construct a rooted tree
$\mathscr T$ on its distinct offsets as follows.  Start at $b_0=0$.
Whenever a step reaches an offset not visited previously, add that
vertex and join it to the current vertex by the step just taken.
Every other step is a \emph{return}.  Write $r_*$ for the number
of returns, and $b(v)$ for the integer offset of a tree vertex $v$.
Thus $|E(\mathscr T)|=\ell-r_*$ and all $b(v)$ are distinct.
Tree edges inherit their labels and signs from their adding steps.
A \emph{tree prime} is a prime dividing a tree-edge label.
\Cref{fig:first-visit-tree} shows a six-step example.

\begin{figure}[htbp]
\centering
\begin{tikzpicture}[>=Stealth,thick,
  vertex/.style={circle,draw,inner sep=2pt},
  every edge/.style={font=\small}]
 \node[vertex] (o) at (0,0) {$0$};
 \node[vertex] (u) at (2.1,0) {$u$};
 \node[vertex] (v) at (4.3,1) {$v$};
 \node[vertex] (w) at (4.3,-1) {$w$};
 \draw[->] (o) to[bend left=12] node[above] {$1$} (u);
 \draw[->,dashed] (u) to[bend left=12] node[below] {$6$} (o);
 \draw[->] (u) to[bend left=12] node[above] {$2$} (v);
 \draw[->,dashed] (v) to[bend left=12] node[pos=.25,below] {$3$} (u);
 \draw[->] (u) to[bend right=12] node[below] {$4$} (w);
 \draw[->,dashed] (w) to[bend right=12] node[pos=.25,above] {$5$} (u);
\end{tikzpicture}
\caption{First visits (solid) and returns (dashed) in the walk
$0\to u\to v\to u\to w\to u\to0$.  Numbers give the chronological
step indices.  The three solid edges form the first-visit tree.
This diagram records only the walk topology.}
\label{fig:first-visit-tree}
\end{figure}

Recall that an occurrence $(i,p)$ is lit when
$p\mid n+b_{i-1}$.  Set
\begin{equation}\label{eq:fixed-line-labels}
 \mathcal F=\{p:\ p\text{ occurs on the line and either }
        p\in\mathcal C\text{ or }p\text{ occurs at least twice}\}.
\end{equation}
Write $n_{\mathcal F}$ for these residue coordinates.  We allow our
restrictions to depend on $n_{\mathcal F}$ and call the primes in
$\mathcal F$ the \emph{fixed labels}.  Here ``fixed'' refers to the
residues: the numerical line already specifies every prime value.
The set includes core primes occurring once on the line, but excludes
background core primes absent from it.  For a fixed tree prime $p$,
define its active vertices by
\[
 V_p=\{v\in V(\mathscr T):p\mid n+b(v)\}.
\]
The connected components of $V_p$ are those of the induced subgraph
on these vertices.  Adjacent active vertices force $p$ to divide
their edge label, since $p\nmid h$ for sufficiently large $B$.

A tree prime $p\notin\mathcal F$ is a center prime occurring on
exactly one step of the entire line.  Let $e_p$ denote its tree edge.
Call it \emph{corrupted} if a tree vertex other than the endpoints
of $e_p$ has an offset congruent to those endpoints modulo $p$.
This is a condition on the numerical line alone.

\begin{proposition}[Retained configurations]\label{prop:retained-lines}
For each closed line there is a predicate
$R_{\boldsymbol d}(n_{\mathcal F})\in\{0,1\}$ with the following
properties.  Whenever it equals one:
\begin{enumerate}
\item all core line occurrences are lit; no primitive specification
with support contained in $\mathcal F$ qualifies at a trace vertex;
and fewer than $t_0$ repeated center occurrences are unlit;
\item every fixed tree prime has $O(B^\rho)$ active components;
fewer than $B^{1-2\rho}$ fixed tree primes have more than one
active component;
\item fewer than $B^{1-2\rho}$ nonfixed tree primes are corrupted.
\end{enumerate}
The implicit constant is uniform in the line.  Moreover,
\begin{equation}\label{eq:retained-line-error}
 \sum_{\boldsymbol d}
  \E\left[|K_{\boldsymbol d}(n)|
       \prod_{j=0}^{\ell}\1_Y(n+b_j)
       (1-R_{\boldsymbol d}(n_{\mathcal F}))\right]
 \ll \exp(-B^{1+\eps/2}).
\end{equation}
The same exceptional estimates used in this proof remain valid
after adding the qualification indicators of at most $2t_0$
specifications and summing their recorded data, for the classes
discarded by an arithmetic saving.  Configurations excluded because
of a fixed forbidden specification instead vanish with the original
$Y$ factors.
\end{proposition}

We prove the proposition in three stages.  First we make active
sets connected on a small number of blocks.  Then a two-block
argument rules out many disconnected fixed labels.  Finally we
handle the accidental offset coincidences of singleton labels.

\subsection{Short gaps and connected blocks}

\begin{lemma}[Hitting paths in a forest]\label{lem:forest-path-hitting}
Let a finite forest carry a finite family of nonempty edge paths.
If the family has no $t$ pairwise edge-disjoint paths, some set of
fewer than $2t$ edges meets every path in the family.
\end{lemma}

\begin{proof}
Root each component.  The top of a path is its vertex nearest that
root.  Choose a path whose top has maximum depth, and mark its one
or two edges incident with the top.  Every path sharing an edge
with the chosen path must contain a marked edge.  Indeed, a path
sharing an edge below a mark and avoiding that mark remains wholly
in the component below it; its top would be strictly deeper.
Remove the paths hit by the marks and repeat.  The paths chosen at
successive stages are pairwise edge-disjoint.  There are fewer than
$t$ stages, so fewer than $2t$ marked edges suffice.
\end{proof}

\begin{lemma}[Connected activity on blocks]\label{lem:activity-blocks}
Except for configurations of zero contribution in
\eqref{eq:retained-line-error} or total absolute contribution
$O(\exp(-B^{1+\eps/2}))$, one can cut $O(t_0)$ edges of the
first-visit tree and cover the resulting forest by $O(B^\rho)$
connected blocks of diameter at most $L$.  The blocks cover every
remaining edge and every vertex of the original tree.  In each block,
the active set of every fixed tree prime is connected or empty.
Consequently each fixed tree prime has $O(B^\rho)$ whole-tree active
components.
\end{lemma}

\begin{proof}
An unlit core occurrence makes $K_{\boldsymbol d}=0$.  A primitive
specification supported on $\mathcal F$ and qualifying at a trace
vertex makes the corresponding $Y$ factor zero for every remaining
coordinate assignment.  Discard these configurations.  By
\Cref{lem:crude-count}, those with at least $t_0$ unlit repeated
center occurrences have the asserted negligible total.  Each of
these decisions uses only the line and $n_{\mathcal F}$.

Cut every tree edge having an unlit fixed occurrence.  There are
fewer than $t_0$ such edges.  In the remaining forest a \emph{gap
for $p$} is a path of length at most $L$ with active endpoints and
at least one interior vertex, all interior vertices being inactive,
for a fixed tree prime $p$.  Such a gap has no $p$-edge: an active
endpoint of a $p$-edge forces its
other endpoint to be active, while a $p$-edge with inactive
endpoints was cut.

Every gap contains a nonfixed prime label.  Otherwise all its
edge labels are fixed and lit, so the path is fully divisible.
Its nonzero total displacement is divisible by $p$, and its edges
avoid $p$.  Taking extra prime $p$ and the whole path as suffix
gives a qualifying specification.  By
\Cref{lem:primitive-descent}, a primitive one with fixed support
qualifies at a vertex of that tree path.  Every tree vertex is a
trace vertex, contrary to the preceding exclusion.

If there are $t_0$ edge-disjoint gaps, choose a nonfixed prime $q$
on each.  Each chosen prime occurs once on the full line, hence
nowhere on the other selected gaps.  The displacement of its gap
is linear in $q$ modulo the corresponding fixed prime $p$, with
coefficient $\pm h$ times the other prime factors of that edge.
This coefficient is invertible modulo $p$, since the gap is
$p$-free and $p\nmid h$.  The selected constraints do not involve
one another's variables.  Their paths and variables have the
description cost of \Cref{lem:triangular-saving}, which discards
this class.  The inequality $t_0\ge B^{2\eps}$ holds for large $B$.

Otherwise \Cref{lem:forest-path-hitting} supplies fewer than
$2t_0$ further cuts meeting every gap.  For each remaining piece
traverse every edge twice and split that traversal into segments
of at most $L$ steps.  The edges and vertices in a segment form
a connected subtree of diameter at most $L$; include a singleton
block for a piece with no edges.  There are
\[
 O(\ell/L+t_0+1)=O(B^\rho)
\]
blocks.  They may overlap.  Any two active vertices in a block
have their joining path inside that block.  An inactive portion
on it would give a surviving short gap, impossible after the cuts.
Thus activity is connected or empty on each block.  Each block
meets at most one whole-tree active component, and every such
component has a vertex in a block.  This proves the component bound.
\end{proof}

\subsection{Comparing two blocks}

The active sets are now connected within each block.  To compare
different components for a fixed prime $p$, suppose blocks $1$ and
$2$ meet those components.  Choose a root $c_j$ in block $j$ and an
active representative $x_p^j$ there.  Since the representatives have
distinct integer offsets,
\begin{equation}\label{eq:two-block-constraint}
 p\mid K+\bigl(b(x_p^2)-b(c_2)\bigr)
           -\bigl(b(x_p^1)-b(c_1)\bigr)\ne0,
 \qquad K=b(c_2)-b(c_1).
\end{equation}
To apply the divisibility case of \Cref{lem:triangular-saving}, we
seek roots and representatives for which both local paths contain
no edge label divisible by $p$.  For many selected primes at once,
we also need a common order in which every selected prime appearing
on either path comes earlier than $p$.  Once $K$ is specified separately, these
conditions give the required triangular dependence.  The next lemma
supplies the choices.

\begin{lemma}[Rerooting two blocks]\label{lem:reroot-constraints}
Let two connected blocks of a tree and $m$ distinct prime indices
$p$ be given.  For each index let $V_p$ be a vertex subset, and let
$A_p^j$ be its intersection with block $j$.  Suppose these
intersections are nonempty and connected.  Edges have labels
containing at most $J$ prime factors.  In either block, an edge has
both endpoints in $A_p^j$ if and only if its label contains $p$.
Assume also that $A_p^1$ and $A_p^2$ lie in distinct components of
the subgraph induced by $V_p$ in the whole tree.

There are roots $c_j$ in block $j$, a set of
$\gg m^{1/8}/(2J+1)$ retained indices, and an order on those indices
with the following property.  If $x_p^j$ is the nearest vertex of
$A_p^j$ to $c_j$, both paths from $c_j$ to $x_p^j$ avoid $p$, and
any retained prime label on either path occurs earlier in the order.
The two representatives $x_p^1,x_p^2$ are distinct.
\end{lemma}

\begin{proof}
Start with arbitrary roots.  A connected vertex set in a tree
contains the entire path between any two of its vertices.  It
therefore has a unique nearest vertex to a chosen root.  The
root-to-$x_p^j$ path avoids $p$-edges.  If a retained label $q$
occurs on this path, both ends of its edge lie in $A_q^j$;
the nearest vertex $x_q^j$ is then a strict ancestor of $x_p^j$.
Thus dependencies are contained in the two strict ancestor orders.
Indices with equal nearest points are incomparable.

Every finite partially ordered set with $u$ elements has a chain
or an antichain of size at least $\sqrt u$: partition its elements
by the length of a longest chain ending there.  Apply this to the
first ancestor order.  If an antichain is obtained, there are no
dependencies from block 1 among its indices, so depth order in
block 2 suffices.  Otherwise apply the same argument in block 2
on the selected chain.  An antichain there similarly suffices.
In the remaining case we have chains in both blocks with at least
$m^{1/4}$ indices.  The Erd\H{o}s--Szekeres monotone-subsequence theorem
\cite[p.~467]{ErdosSzekeres1935} shows that two total orders have a
common agreeing or reversed subsequence of size at least the square
root of their length.  To recall the elementary argument, assign to each position the lengths
of the longest increasing and decreasing subsequences ending
there.  Two positions have different pairs, so if both lengths
are less than $\sqrt u$ there cannot be $u$ positions.  We thus
retain at least a constant times $m^{1/8}$ indices.  Agreeing
orders already give the conclusion.

Suppose the orders are reversed.  In block 2 take the segment
from its root to the farthest selected $x_p^2$.  The intersection
of each $A_p^2$ with that segment is a closed vertex interval
beginning at $x_p^2$; these left endpoints are distinct.  At a
vertex of the segment, each positive-length interval containing
it contains an incident segment edge.  Each of the at most two
such edges has at most $J$ prime factors.  At most one interval
can consist of that vertex alone.  Hence at most $2J+1$ intervals
overlap at a vertex.  Greedily coloring closed intervals in order
of their left endpoints uses at most $2J+1$ colors: an existing
color is available only if its preceding right endpoint is
strictly earlier than the new left endpoint.  One color retains
at least a $1/(2J+1)$ fraction of pairwise vertex-disjoint intervals.

Reroot block 2 at the far end of the segment.  The nearest point
of $A_p^2$ is now the right endpoint of its interval.  Indeed, for
any vertex of $A_p^2$ off the segment, its projection onto the
segment lies in that interval, because the path from $x_p^2$ to
the vertex lies in $A_p^2$.  The path from the new root first
reaches $A_p^2$ at the right endpoint.  This also proves that
off-segment branches cannot change the nearest point.

The right endpoints of vertex-disjoint intervals have the same
order as their left endpoints along the original segment.  Seen
from its other end, that order reverses.  The two block orders
therefore agree after rerooting.  The first paragraph applies to
the new root as well and proves the dependency assertion.  Each
representative stayed in the same whole-tree component for its
prime, so the two representatives for that prime are still distinct.
\end{proof}

\begin{lemma}[Few disconnected fixed labels]\label{lem:few-disconnected}
Among the configurations retained by \Cref{lem:activity-blocks},
those with at least $B^{1-2\rho}$ fixed tree primes having multiple
active components have total absolute contribution
$O(\exp(-B^{1+\eps/2}))$.
\end{lemma}

\begin{proof}
For every such prime choose two blocks meeting distinct whole-tree
active components.  There are $O(B^{2\rho})$ ordered block pairs,
so one pair serves $m\gg B^{1-4\rho}$ primes.  Inside each block,
the active set is nonempty and connected.  Every edge carrying a
fixed prime has active endpoints after the earlier cuts, and the
converse follows from $p\nmid h$.  Apply
\Cref{lem:reroot-constraints} and then specify the common integer
\[
 K=b(c_2)-b(c_1),\qquad |K|=O(\ell H).
\]
Each retained prime gives the test \eqref{eq:two-block-constraint}.
The expression is independent of $p$, since both local paths
avoid it.  It involves only earlier retained variables by the
lemma.  Its nonvanishing follows from distinctness of the two
representatives and the distinct integer offsets of tree vertices.
The number of tests is
\[
 \gg \frac{B^{(1-4\rho)/8}}{J}
   =\frac{B^{1/10}}{O(\log B)}>B^{2\eps}
\]
for large $B$.

Roots, endpoints, and prime choices can be recorded with the
entropy allowed in \Cref{lem:triangular-saving}.  Choose $K$ only
after the rerooting; its $O(\ell H+1)$ possibilities cost
$\exp(O(B))$.  Although its original value depends on the labels,
fixing it and dropping that defining equation enlarges a positive
sum.  The local path expressions then have the required triangular
dependence.  Retain their nonzero tests under this enlargement and
apply \Cref{lem:triangular-saving}.
\end{proof}

\subsection{Accidental coincidences of singleton labels}

The preceding argument concerned fixed residue coordinates.  A
nonfixed tree prime requires a different test: its residue must
remain free, so we use only numerical coincidences between offsets.

\begin{lemma}[Few corrupted singleton labels]\label{lem:few-corrupted}
Lines having at least $B^{1-2\rho}$ corrupted nonfixed tree primes
have total absolute contribution
$O(\exp(-B^{1+\eps/2}))$.
\end{lemma}

\begin{proof}
For every corrupted prime choose a witnessing vertex outside its
edge endpoints, and put $D=\lfloor B^{1-3\rho}\rfloor$.  Distance
from a vertex to an edge means its minimum tree distance to the two
endpoints.  Suppose first that at least half the selected witnesses have distance at
most $D$ from their prime edge.  Use the path from the nearer edge
endpoint to its witness.  This path avoids the prime's sole edge,
so its nonzero displacement is independent of that prime and is
divisible by it.  The expression contains at most $JD$ other
prime labels.

On the selected primes draw a directed dependency edge $p\to q$
if the test for $p$ contains $q$.  Its outdegree is at most $JD$.
Every induced subgraph of the underlying undirected graph has
average degree at most $2JD$.  Successively choose a vertex of
degree at most $2JD$ and delete its neighbors.  The resulting
independent set has size at least a constant times
\[
 \frac{B^{1-2\rho}}{JD+1}\gg\frac{B^\rho}{J}>B^{2\eps}.
\]
Its tests have no selected-variable dependencies.  They are
covered by \Cref{lem:triangular-saving}.

Otherwise at least half the witnesses have distance greater than
$D$.  Cover the full tree by $O(\ell/D+1)=O(B^{3\rho})$ connected
traversal blocks of diameter at most $D$, as in the earlier
construction.  Assign each prime edge to a block containing that
edge and its witness to a block containing that vertex.  One
ordered pair serves $\gg B^{1-8\rho}$ primes.  No selected prime
edge has an endpoint in the second block: its corresponding
witness lies there, so that would put it at distance at most $D$.
Consequently paths inside the second block contain no selected
prime label.

Root the two blocks at $c_1,c_2$.  In block 1 let $x_p^1$ be the
nearer endpoint of the prime edge; in block 2 let $x_p^2$ be its
witness.  The first local path avoids $p$.  If it contains a
selected prime $q$, the nearer endpoint of the $q$-edge is a
strict ancestor of $x_p^1$.  Thus the tests
\eqref{eq:two-block-constraint}, after specifying the common root
offset, are triangular in the depths of the first-block
endpoints.  The second local paths introduce no selected
variables.  Each tested displacement is nonzero because its
witness differs from the chosen edge endpoint.  Since
$1-8\rho=3/5>2\eps$, there are more than enough tests for
\Cref{lem:triangular-saving}.  All choices of witnesses, blocks,
roots, and paths fit its description bound.
\end{proof}

\begin{proof}[Proof of \Cref{prop:retained-lines}]
Define $R_{\boldsymbol d}$ by retaining the conditions obtained
in \Cref{lem:activity-blocks,lem:few-disconnected,lem:few-corrupted}.
Every activity and occurrence status of a fixed label is determined
by $n_{\mathcal F}$; fixed forbidden qualifications use only those
coordinates; corruption uses only numerical offsets.  Covers and
witnesses can be chosen by fixed finite ordering whenever they
exist.  Thus $R_{\boldsymbol d}$ depends on no free center residue.

The zero exclusions vanish with $K_{\boldsymbol d}\prod_j\1_Y$.
The other excluded classes have the stated absolute contribution
by the three lemmas, yielding \eqref{eq:retained-line-error}.
Their proofs used precisely the majorants of
\Cref{lem:crude-count,lem:triangular-saving}, which allow up to
$2t_0$ attached specifications.  Their retained tests depend only
on the numerical line and fixed residues, and are unchanged when
qualification indicators are added.  This proves the additional assertion as well.
\end{proof}

We may therefore impose the retained predicate before estimating
the remaining conditional averages.  This preserves the residue
coordinates of singleton center labels.  The next section deals
with the remaining dependence of the factors $\1_Y$ on those
coordinates by expanding them into short lists of qualification
conditions.
